% ============================================================
%  Human Detection — Algorithm Evaluation Report
%  Smart Thermal System for Patient Safety Monitoring
%  Afeka College of Engineering, Tel-Aviv
%
%  Compiled with: pdflatex (TeX Live / Overleaf)
% ============================================================
\documentclass[11pt, a4paper]{article}

% ---- Packages -----------------------------------------------
\usepackage[a4paper, margin=2.5cm]{geometry}
\usepackage{booktabs}
\usepackage{multirow}
\usepackage{array}
\usepackage{amsmath}
\usepackage{amssymb}
\usepackage[table]{xcolor}   % loads colortbl for \cellcolor
\usepackage{graphicx}
\usepackage{caption}
\usepackage{hyperref}
\usepackage{parskip}
\usepackage{enumitem}
\usepackage{microtype}
\usepackage{float}

% ---- Colours ------------------------------------------------
\definecolor{tpgreen}{RGB}{0,128,0}
\definecolor{tnblue}{RGB}{31,73,125}
\definecolor{fpred}{RGB}{192,0,0}
\definecolor{fnyellow}{RGB}{150,100,0}
\definecolor{headerblue}{RGB}{31,73,125}
\definecolor{cellgreenlight}{RGB}{198,239,206}
\definecolor{cellbluelight}{RGB}{221,235,247}
\definecolor{cellredlight}{RGB}{255,199,206}
\definecolor{cellyellowlight}{RGB}{255,235,156}

% ---- Hyperref setup -----------------------------------------
\hypersetup{
    colorlinks = true,
    linkcolor  = headerblue,
    urlcolor   = headerblue,
    citecolor  = headerblue
}

% ---- Confusion matrix macro ---------------------------------
% Usage: \confmat{TP}{FN}{FP}{TN}{pos_total}{neg_total}
% \def saves args before entering tabular (avoids # conflicts).
\newcommand{\confmat}[6]{%
    \def\cmTP{#1}\def\cmFN{#2}\def\cmFP{#3}\def\cmTN{#4}%
    \def\cmPos{#5}\def\cmNeg{#6}%
    \vspace{4pt}
    \begin{center}
    \renewcommand{\arraystretch}{1.8}%
    \begin{tabular}{lcc}
        \hline
        \rule{0pt}{2.4ex}%
        & \textbf{Pred:\ Human} & \textbf{Pred:\ Empty} \\[2pt]
        \hline
        \rule{0pt}{2.6ex}%
        \textbf{Truth:\ Human} (\cmPos\ frames)
            & \cellcolor{cellgreenlight}\textcolor{tpgreen}{\textbf{TP}}\;$=$\;\cmTP
            & \cellcolor{cellyellowlight}\textcolor{fnyellow}{\textbf{FN}}\;$=$\;\cmFN \\[4pt]
        \textbf{Truth:\ Empty} (\cmNeg\ frames)
            & \cellcolor{cellredlight}\textcolor{fpred}{\textbf{FP}}\;$=$\;\cmFP
            & \cellcolor{cellbluelight}\textcolor{tnblue}{\textbf{TN}}\;$=$\;\cmTN \\[2pt]
        \hline
    \end{tabular}
    \end{center}
    \vspace{4pt}
}

% ---- Title --------------------------------------------------
\title{%
    {\large Afeka College of Engineering, Tel-Aviv}\\[0.4em]
    {\large Smart Thermal System for Patient Safety Monitoring}\\[0.8em]
    {\LARGE\bfseries Human Detection — Algorithm Evaluation Report}\\[0.4em]
    {\normalsize MLX90640 Thermal Sensor \quad 22 Scenes \quad
     80\,/\,10\,/\,10 Train\,/\,Val\,/\,Test Split}
}
\author{%
    Guy Chen \and Yaniv Blau \and Roy Lieberman\\[0.3em]
    \small Supervisor: Or Zilberberg \qquad Advisor: Dr.\@ Oshrit Hoffer
}
\date{June 2026}

% =============================================================
\begin{document}
\maketitle
\tableofcontents
\newpage

% =============================================================
\section{Overview}
% =============================================================

This report presents a systematic evaluation of four human detection algorithms applied
to thermal imagery from the MLX90640 32$\times$24 sensor.  Two evaluation criteria are
applied to each algorithm: a lenient \emph{frame-level presence} criterion and a strict
\emph{localisation} criterion requiring IoU\,$>\,0.5$ between predicted and ground-truth
bounding boxes.  The goal is to characterise which algorithms are suitable for each
downstream task in the safety monitoring pipeline.

% =============================================================
\section{Dataset and Experimental Setup}
% =============================================================

\subsection{Dataset}

Recordings were collected at the Beer Sheva Mental Health Center using three MLX90640
sensors mounted in the ward.  Each sensor produces 32$\times$24 thermal frames at 8\,Hz.
The evaluation uses \textbf{22 scenes} (two scenes — \texttt{2pplrun} and
\texttt{3pplfight} — were excluded pending annotation review).  Scene names describe the
scenario: single or multiple persons, presence of a heat source, running, falling, etc.

Preprocessing is applied uniformly across all algorithms: each channel is passed through
a \textbf{Tateno pipeline} (Gaussian denoising $\to$ background subtraction $\to$ L1
rectification), calibrated on even-indexed frames from the empty-room recording.

\subsection{Train\,/\,Validation\,/\,Test Split}

For each \emph{(scene, channel)} group, labeled frame indices are sorted by temporal
order and split sequentially: the first 80\,\% form the training set, the next 10\,\%
the validation set, and the final 10\,\% the test set.  Sequential splitting prevents
temporal leakage while ensuring every scene contributes proportionally to every split.

\begin{table}[H]
\centering
\caption{Split statistics (frame-channel examples).}
\renewcommand{\arraystretch}{1.3}
\begin{tabular}{lccc}
\toprule
\textbf{Split} & \textbf{Total}
    & \textbf{Person-present (Positive)}
    & \textbf{Empty (Negative)} \\
\midrule
Train      & 4{,}632 & 4{,}393 & 239 \\
Validation &    576  &    559  &  17 \\
Test       &    621  &    608  &  13 \\
\bottomrule
\end{tabular}
\label{tab:split}
\end{table}

The dataset is strongly imbalanced ($\approx 98$\,\% positive), reflecting real-world
conditions in a monitored ward where persons are present nearly continuously.

% =============================================================
\section{Metric Definitions}
% =============================================================

\subsection{Confusion Matrix}

Every frame is reduced to a binary verdict: \textbf{Human Detected (Positive)} or
\textbf{Empty (Negative)}.  Four outcomes are possible:

\begin{itemize}[leftmargin=2em]
    \item \textbf{TP} (True Positive): person present and correctly detected.
    \item \textbf{TN} (True Negative): no person present and system correctly silent.
    \item \textbf{FP} (False Positive): no person present but alarm raised
          (\emph{false alarm}).
    \item \textbf{FN} (False Negative): person present but missed by the system
          (\emph{missed detection}).
\end{itemize}

\subsection{Derived Metrics}

\begin{align}
    \text{Accuracy}  &= \frac{TP + TN}{TP + TN + FP + FN} \label{eq:acc}  \\[6pt]
    \text{Precision} &= \frac{TP}{TP + FP}                \label{eq:prec} \\[6pt]
    \text{Recall}    &= \frac{TP}{TP + FN}                \label{eq:rec}  \\[6pt]
    \text{F1}        &= \frac{2 \cdot \text{Precision} \cdot \text{Recall}}
                              {\text{Precision} + \text{Recall}}            \label{eq:f1}
\end{align}

\paragraph{Accuracy} (Eq.\,\ref{eq:acc}) is the overall fraction of correctly classified
frames.  On a heavily imbalanced dataset, accuracy is susceptible to optimistic inflation:
a detector that \emph{always fires} achieves $\approx 98$\,\% accuracy while missing
every empty frame.

\paragraph{Precision} (Eq.\,\ref{eq:prec}) is the fraction of predicted positives that
are genuinely positive.  High precision implies few false alarms.  Excessive false
positives cause \emph{alert fatigue}, leading staff to disable or ignore the system.

\paragraph{Recall} (Eq.\,\ref{eq:rec}) is the fraction of genuine positives that are
detected.  High recall implies few missed detections.  Recall is the
\textbf{primary safety metric}: a missed intruder, fallen patient, or fire event is a
direct system failure.

\paragraph{F1-Score} (Eq.\,\ref{eq:f1}) is the harmonic mean of precision and recall.
It balances both objectives and is used as the headline metric because of the class
imbalance.

\subsection{Success Criteria}

Two criteria are applied throughout this report:

\begin{description}[leftmargin=2.5em, style=nextline]
    \item[Criterion A — Frame-level presence (no IoU):]
        The detector is credited with a TP whenever it outputs \emph{any} bounding box
        on a person-present frame, regardless of spatial accuracy.  This criterion
        answers: \emph{``Did the system know someone was in the room?''}  It is
        appropriate for simple presence-only alarming (e.g.\ restricted-area breach
        detection).

    \item[Criterion B — Localisation (IoU\,$>\,0.5$):]
        A frame is counted as TP only if at least one predicted box overlaps a
        ground-truth box by more than 50\,\% Intersection-over-Union (IoU).  This
        criterion answers: \emph{``Did the system correctly locate the person?''}
        Accurate localisation is a prerequisite for passing bounding boxes downstream
        to contact detection.
\end{description}

% =============================================================
\section{Algorithms}
% =============================================================

\begin{table}[H]
\centering
\caption{Summary of evaluated algorithms.}
\renewcommand{\arraystretch}{1.3}
\begin{tabular}{>{\raggedright}p{2.9cm}
                >{\raggedright}p{2.4cm}
                >{\raggedright}p{2.0cm}
                >{\raggedright\arraybackslash}p{6.8cm}}
\toprule
\textbf{Algorithm} & \textbf{Type} & \textbf{Training} & \textbf{Description} \\
\midrule
Adaptive Threshold
    & Rule-based CV
    & None
    & Adaptive Gaussian thresholding, morphological closing, and geometric blob
      filtering (area, solidity, aspect ratio). \\[4pt]
HOG\,+\,SVM
    & Classical ML
    & LinearSVC
    & Sliding window ($14\!\times\!6$\,px) with Histogram of Oriented Gradients
      features; Linear SVM classifier; single scale. \\[4pt]
HOG\,+\,SVM\,+\,Pyramid
    & Classical ML
    & Same as above
    & Identical trained SVM weights; inference at scales
      $\{0.75\times,\,1.00\times,\,1.25\times,\,1.50\times\}$;
      unified NMS across all scales. \\[4pt]
MobileNet-SSD
    & Deep learning
    & PyTorch (CUDA)
    & Micro-MobileNet backbone with SSD anchor regression heads;
      30 training epochs; RTX~2060~Super. \\
\bottomrule
\end{tabular}
\label{tab:algorithms}
\end{table}

% =============================================================
\section{Results — Criterion A: Frame-Level Presence}
% =============================================================

Under Criterion~A, a detection anywhere on a person-present frame is counted as correct.

\subsection{Confusion Matrices}

\subsubsection*{Adaptive Threshold}
\confmat{591}{17}{1}{12}{608}{13}

\subsubsection*{HOG\,+\,SVM (single scale)}
\confmat{585}{23}{3}{10}{608}{13}

\subsubsection*{MobileNet-SSD}
\confmat{593}{15}{0}{13}{608}{13}

\subsection{Metric Summary}

\begin{table}[H]
\centering
\caption{Criterion A — test-set metrics (frame-level presence, no IoU requirement).}
\renewcommand{\arraystretch}{1.3}
\begin{tabular}{lcccccccc}
\toprule
\textbf{Algorithm}
    & \textbf{Acc} & \textbf{Prec} & \textbf{Rec} & \textbf{F1}
    & \textbf{TP}  & \textbf{TN}   & \textbf{FP}  & \textbf{FN} \\
\midrule
Adaptive Threshold
    & 97.1\% & 99.8\%          & 97.2\% & \textbf{98.5\%}
    & 591    & 12              &  1     & 17 \\
HOG\,+\,SVM
    & 95.8\% & 99.5\%          & 96.2\% & \textbf{97.8\%}
    & 585    & 10              &  3     & 23 \\
MobileNet-SSD
    & 97.6\% & \textbf{100.0\%}& \textbf{97.5\%} & \textbf{98.8\%}
    & 593    & 13              &  0     & 15 \\
\bottomrule
\end{tabular}
\label{tab:crA}
\end{table}

\subsection{Interpretation}

Under the lenient criterion, \textbf{all three algorithms perform well}, with F1-scores
between 97.8\,\% and 98.8\,\%.  The results demonstrate that \emph{detecting the
presence} of a person in a thermal frame is a tractable problem even for a
parameter-free, rule-based approach.

\textbf{MobileNet-SSD} leads at F1\,=\,98.8\,\% and is the only algorithm to produce
\textbf{zero false positives} — it never raised an alarm on an empty room.  Its 15 missed
detections are mostly frames with partial occlusion or atypical thermal profiles.

\textbf{Adaptive Threshold} is effectively tied (F1\,=\,98.5\,\%) with no training cost.
Its single false positive originates from a thermal reflection artefact in the empty-room
recording.

\textbf{HOG\,+\,SVM} ranks third at F1\,=\,97.8\,\%, with 3 false alarms: the sliding
window found background thermal gradients that matched its learnt person template.

The conclusion from Criterion~A alone would be misleadingly optimistic.  The next section
exposes the true differences between algorithms.

% =============================================================
\section{Results — Criterion B: Localisation (IoU\,$>\,0.5$)}
% =============================================================

Under Criterion~B, a frame only counts as a TP if the predicted bounding box overlaps the
ground-truth annotation by IoU\,$>\,0.5$.  Frames where a person was detected but the box
was misplaced become false negatives.

\subsection{Without Scale Pyramid}

\subsubsection*{Adaptive Threshold}
\confmat{423}{185}{1}{12}{608}{13}

\subsubsection*{HOG\,+\,SVM (single scale)}
\confmat{246}{362}{3}{10}{608}{13}

\subsubsection*{MobileNet-SSD}
\confmat{549}{59}{0}{13}{608}{13}

\subsection{With Scale Pyramid (HOG\,+\,SVM)}

The scale pyramid runs the same trained SVM weights on four rescaled copies of each frame
($0.75\times$, $1.00\times$, $1.25\times$, $1.50\times$), projects all detections back
to original coordinates, and performs a single unified NMS pass.

\subsubsection*{HOG\,+\,SVM\,+\,Pyramid}
\confmat{376}{232}{12}{1}{608}{13}

\subsection{Metric Summary}

\begin{table}[H]
\centering
\caption{Criterion B — test-set metrics (IoU\,$>\,0.5$ localisation requirement).}
\renewcommand{\arraystretch}{1.3}
\begin{tabular}{lcccccccc}
\toprule
\textbf{Algorithm}
    & \textbf{Acc} & \textbf{Prec} & \textbf{Rec} & \textbf{F1}
    & \textbf{TP}  & \textbf{TN}   & \textbf{FP}  & \textbf{FN} \\
\midrule
Adaptive Threshold
    & 70.0\% & 99.8\%          & 69.6\% & \textbf{82.0\%}
    & 423    & 12 &  1 & 185 \\
HOG\,+\,SVM
    & 41.2\% & 98.8\%          & 40.5\% & \textbf{57.4\%}
    & 246    & 10 &  3 & 362 \\
HOG\,+\,SVM\,+\,Pyramid
    & 60.7\% & 96.9\%          & 61.8\% & \textbf{75.5\%}
    & 376    &  1 & 12 & 232 \\
MobileNet-SSD
    & 90.5\% & \textbf{100.0\%}& \textbf{90.3\%} & \textbf{94.9\%}
    & 549    & 13 &  0 &  59 \\
\bottomrule
\end{tabular}
\label{tab:crB}
\end{table}

\subsection{Interpretation}

The localisation criterion separates the algorithms dramatically: nearly all new false
negatives are frames where the system \emph{did} detect a person but placed the box
inaccurately.

\paragraph{MobileNet-SSD} remains the clear leader at F1\,=\,94.9\,\%
($-3.9$\,pp from Criterion~A).  Its SSD anchor regression heads are explicitly trained
to minimise bounding-box error, making it the only algorithm that genuinely
\emph{localises} persons rather than merely detecting their presence.  The 59 false
negatives are concentrated in scenes with atypical postures (falling, crouching) and
dense multi-person occlusion.  Crucially, \textbf{zero false alarms} are produced on
empty-room frames.

\paragraph{Adaptive Threshold} drops to F1\,=\,82.0\,\% ($-16.5$\,pp).  The
morphological closing reliably finds the hottest blob in a frame, but the resulting
contour bounding rectangle is seldom tight enough to clear IoU\,=\,0.5 against a manually
drawn annotation.  The 185 false negatives are almost entirely frames where a person
\emph{was} detected but the blob shape diverged from the ground-truth box.

\paragraph{HOG\,+\,SVM (single scale)} collapses to F1\,=\,57.4\,\%
($-40.4$\,pp) — the worst localisation result.  The root cause is the \textbf{fixed
sliding window} ($14\!\times\!6$\,px on a $32\!\times\!24$ sensor).  When a person
appears at a scale or aspect ratio different from the training window, the best achievable
IoU is bounded well below 0.5 regardless of SVM confidence: detection occurs, but
localisation does not.

\paragraph{HOG\,+\,SVM\,+\,Pyramid} partially recovers to F1\,=\,75.5\,\%
($+18.1$\,pp over single scale).  Running inference at four scales allows the fixed window
to match persons across a range of apparent distances, converting 130 additional false
negatives to true positives.  However, the pyramid also introduces a false-alarm
regression: FP rises from~3 to~12, and true negatives on empty frames fall from~10 to~1.
The multi-scale sweep provides more recall on positive frames at the cost of substantially
reduced selectivity on empty backgrounds.

% =============================================================
\section{Consolidated Comparison}
% =============================================================

\begin{table}[H]
\centering
\caption{Complete metric comparison — test set, all algorithms, both criteria.}
\renewcommand{\arraystretch}{1.35}
\begin{tabular}{llcccc}
\toprule
\textbf{Algorithm} & \textbf{Criterion}
    & \textbf{Accuracy} & \textbf{Precision} & \textbf{Recall} & \textbf{F1} \\
\midrule
\multirow{2}{*}{Adaptive Threshold}
    & No IoU          & 97.1\% & 99.8\%          & 97.2\% & 98.5\% \\
    & IoU\,$>\,0.5$   & 70.0\% & 99.8\%          & 69.6\% & 82.0\% \\
\addlinespace
\multirow{2}{*}{HOG\,+\,SVM}
    & No IoU          & 95.8\% & 99.5\%          & 96.2\% & 97.8\% \\
    & IoU\,$>\,0.5$   & 41.2\% & 98.8\%          & 40.5\% & 57.4\% \\
\addlinespace
HOG\,+\,SVM\,+\,Pyramid
    & IoU\,$>\,0.5$   & 60.7\% & 96.9\%          & 61.8\% & 75.5\% \\
\addlinespace
\multirow{2}{*}{MobileNet-SSD}
    & No IoU          & 97.6\% & \textbf{100.0\%}& 97.5\% & 98.8\% \\
    & IoU\,$>\,0.5$   & 90.5\% & \textbf{100.0\%}& \textbf{90.3\%} & \textbf{94.9\%} \\
\bottomrule
\end{tabular}
\label{tab:full}
\end{table}

The F1 drop from Criterion~A to Criterion~B quantifies each algorithm's
\emph{localisation gap}:

\begin{table}[H]
\centering
\caption{Localisation gap: F1 under Criterion~A vs.\ Criterion~B.}
\renewcommand{\arraystretch}{1.3}
\begin{tabular}{lccc}
\toprule
\textbf{Algorithm}
    & \textbf{F1 (no IoU)} & \textbf{F1 (IoU\,$>\,0.5$)} & \textbf{Drop} \\
\midrule
MobileNet-SSD          & 98.8\% & 94.9\% & $-3.9$\,pp \\
Adaptive Threshold     & 98.5\% & 82.0\% & $-16.5$\,pp \\
HOG\,+\,SVM            & 97.8\% & 57.4\% & $-40.4$\,pp \\
HOG\,+\,SVM\,+\,Pyramid& ---    & 75.5\% & $-22.3$\,pp vs.\ no-IoU \\
\bottomrule
\end{tabular}
\label{tab:gap}
\end{table}

% =============================================================
\section{Conclusions}
% =============================================================

\subsection{Presence Detection (Criterion A)}

All three base algorithms achieve F1\,$\geq$\,97.8\,\% under the frame-level presence
criterion.  For applications requiring only a binary room-occupancy signal (e.g.\
restricted-area breach detection), even the parameter-free Adaptive Threshold detector is
sufficient, with the advantage of requiring no training data or GPU resources.

\subsection{Localisation (Criterion B)}

Only \textbf{MobileNet-SSD} meets a practical localisation threshold, achieving
F1\,=\,94.9\,\% with zero false alarms.  It is the recommended algorithm for any
downstream task that consumes bounding boxes — most importantly, contact detection, where
the foot-point projection from each person's bounding box is the input to the multi-view
fusion stage.

Adaptive Threshold achieves F1\,=\,82.0\,\% and may be acceptable as a lightweight
fallback, but its 30\,\% false-negative rate on localisation means that roughly one in
three person-present frames will not produce a usable bounding box.

HOG\,+\,SVM, even with scale pyramid (F1\,=\,75.5\,\%), remains 19\,pp below
MobileNet-SSD.  The pyramid improves recall by $+18$\,pp but introduces an unacceptable
false-alarm rate on empty frames (92\,\% false-alarm rate on the empty-room test split).
Further improvement would require \textbf{hard negative mining} to suppress background
misfires, followed by a \textbf{bounding-box regression layer} to correct residual
window-position errors.

\subsection{Recommendation}

\begin{itemize}
    \item \textbf{Primary detector:} MobileNet-SSD — best F1, zero false positives,
          suitable for bounding-box–dependent downstream tasks.
    \item \textbf{Lightweight fallback:} Adaptive Threshold — no training required,
          adequate for presence-only alarming.
    \item \textbf{HOG\,+\,SVM:} Not recommended for localisation tasks in its current
          form; viable only as a presence-only detector (without pyramid).
\end{itemize}

\end{document}
